% ===================================================================
%  TÁBLA Uimh. 8 — An Fómhar. — An tSealgaireacht, Baint na bhFíonchaor,
%  An Iascaireacht.
%  Traduction 2026 d'apres texto/io, controlee sur texto/fr, texto/en
%  et texto/es. Consigne : voir texto/ga/10-tabla-01.tex.
%
%  L'appel de note est dans le TITRE de la scene — « An Fómhar (*) ».
%  LE BLOC c2-03-1 EST UN ECART DECLARE dans renvoji.py.
%
%  LE TABLEAU DE LA MOISSON PROLONGE CELUI DE LA FERME : « corrán » la
%  faucille, « speal » la faux, « arbhar » le ble, « dias » l'epi,
%  « punann » la gerbe — tous celtiques anciens, comme le troupeau du
%  tableau 7.
% ===================================================================

%%K t08-noto-1 noto
\VUnotes{143.2mm}{%
\textit{(*) Toisc nach bhfuil an focal seo cruinn: fómhar an lín, baint na
bhfíonchaor, baint na n-úll? B'fhéidir gurbh fhearr «} \VUgras{mesono}
\textit{» a ghlacadh, a moladh in} Mondo \textit{XIV, lch. 242. Ach go dtí seo níor
ghlac an tAcadamh leis an bhfréamh nua seo agus tá sí ag fanacht le plé de réir
gnáthrialacha ghluaiseacht an ido.}%
}

%%K t08-apar-1 apar
\VUsaut{5.0mm}
\VUtitre{15.0pt}{60}{TÁBLA Uimh. 8}
\VUsaut{3.0mm}
\VUcentre{13.2pt}{90}{\textit{An Chéad Radharc.}}
\VUsaut{3.0mm}
\VUpk{10.2pt}{40}{An Fómhar (*).}
\VUsaut{4.0mm}
\VUpk{10.2pt}{40}{Aghaidh na bpáirceanna.}
\VUsaut{3.0mm}

%%K t08-c1-01-1 p
1. --- Le teacht an \VUgras{tsamhraidh} tá aghaidh nua ar na páirceanna arís. Tá an
síol a cuireadh san eitre tar éis fás; d'ardaigh planda beag glas as an talamh agus
thug sé a dhias. Líonadh an gráinne, ansin aibigh sé, agus anois níl fágtha ach an
fómhar a bhaint.

%%K t08-c1-02-1 p
2. --- Tá \VUgras{bláthanna na bpáirceanna} ar oscailt. Cuireann na
\VUgras{gormáin} \textsuperscript{(1)}, na \VUgras{poipíní}
\textsuperscript{(2)} agus na \VUgras{méiríní sí} \textsuperscript{(3)} caismirt
mheidhreach i dtón cothrom an arbhair lena \VUgras{gcorólla} úra.

%%K t08-c1-03-1 p
3. --- Tá na crainn torthaí gléasta i nduilliúr; tá na torthaí aibí. Cromann an
\VUgras{crann silíní} beag \textsuperscript{(4)} faoi na \VUgras{silíní} áille
\textsuperscript{(5)}, a bhaineann agus a itheann na páistí le pléisiúr.

%%K t08-c1-04-1 p
4. --- Is féidir leis an tréad bheith saor an lá ar fad anois.

%%K t08-c1-04-2 p
Tá na \VUgras{huain} \textsuperscript{(6)} tar éis fás agus tá siad ina
\VUgras{gcaoirigh} bhreátha \textsuperscript{(7)} agus ina \VUgras{reithí}
bhreátha \textsuperscript{(8)} le holann thiubh. Bíonn siad ag ithe an
\VUgras{fhéir} go socair ar an \VUgras{bhféarach} \textsuperscript{(9)}, atá
breactha le \VUgras{nóiníní}.

%%K t08-c1-04-3 p
Bearrfar an \VUgras{olann} de na hainmhithe seo gan mhoill, olann as a ndéantar an
t-éadach dár gcuid éadaigh.

%%K t08-tit-2 sub
\VUsaut{5.0mm}
\VUpk{10.2pt}{40}{An tAoire.}
\VUsaut{3.4mm}

%%K t08-c1-05-1 p
5. --- Dealraíonn nach mbreathnaíonn an \VUgras{t-aoire} \textsuperscript{(10)}
mórán orthu. Níl a shuim ach sa stoirm atá ag dul thar bráid ag bun na spéire, agus
dealraíonn an \VUgras{buachaill caorach} \textsuperscript{(13)}, ag tabhairt a
\VUgras{bhuidéil} \textsuperscript{(14)} chun a bhéil, níos neamhairdiúla fós.

%%K t08-c1-06-1 p
6. --- Tá an teas ag brú; seo an \VUgras{madra} \textsuperscript{(15)} ag dul chun
é féin a fhuarú; ólann sé uisce úr an \VUgras{tsrutháin} \textsuperscript{(16)} agus
scanraíonn sé \VUgras{frog} beag bocht \textsuperscript{(17)}, a bhí i bhfolach san
fhéar ar an mbruach. Léimeann seisean isteach san uisce glan, mar a bhfuil na
\VUgras{duilleoga báite} \textsuperscript{(18)} faoi bhláth.

%%K t08-c1-07-1 p
7. --- Tá an \VUgras{glasóg} \textsuperscript{(19)} á folcadh féin in aice leis an
\VUgras{tobar} beag \textsuperscript{(20)}, a bhrúchtann idir dhá chloch agus a
théann i bhfolach faoi na \VUgras{sceacha deilgneacha} \textsuperscript{(21)} agus
faoi na \VUgras{toim sceiche gile} \textsuperscript{(22)}.

%%K t08-tit-3 sub
\VUsaut{5.0mm}
\VUpk{10.2pt}{40}{An Fómhar.}
\VUsaut{3.4mm}

%%K t08-c1-08-1 p
8. --- Do pháistí an tsráidbhaile is am an-mheidhreach é baint an arbhair. Bíonn
siad ag \VUgras{tumadh} go meidhreach \textsuperscript{(24)} ar na
\VUgras{cruacha tuí} \textsuperscript{(23)}, cabhraíonn siad leis na
\VUgras{buanaithe} \textsuperscript{(26)}, agus fad is atá siad sin ag baint an
\VUgras{ghoirt} \textsuperscript{(27)} lena \VUgras{speala}
\textsuperscript{(25)}, atá faobhraithe go maith ar an \VUgras{gcloch fhaobhair}
\textsuperscript{(30)}, cruinníonn siadsan an t-arbhar agus ceanglaíonn siad ina
\VUgras{phunanna} \textsuperscript{(31)} nó ina \VUgras{cheanglacha}
\textsuperscript{(32)} é.

%%K t08-c1-09-1 p
9. --- Uaireanta ligtear do na cinn is mó an \VUgras{corrán}
\textsuperscript{(12)} a úsáid chun na \VUgras{gais óir}
\textsuperscript{(28)} a ghearradh, gais a iompraíonn \VUgras{diasa} troma
\textsuperscript{(29)}, lán de ghráinní; agus na cinn bheaga, agus iad ag faire ar
an \VUgras{eallach} \textsuperscript{(33)} atá ag innilt ar an
\VUgras{móinéar} \textsuperscript{(34)} faoi scáth na \VUgras{teile} móire
\textsuperscript{(35)}, beireann siad lena \VUgras{líon} \textsuperscript{(36)} ar
\VUgras{fhéileacáin} áille.

%%K t08-c1-09-2 p
Dreapann cuid acu fiú ar an \VUgras{gcairt} \textsuperscript{(37)}, chun na punanna
a leagan a thugtar dóibh ar an \VUgras{bpíce} \textsuperscript{(38)}.

%%K t08-c1-10-1 p
10. --- Ar fud an \VUgras{mhachaire} ar fad \textsuperscript{(39)}, mar a n-éiríonn
na \VUgras{fuiseoga} \textsuperscript{(41)}, tá fuadar mór. Tá gach duine gafa leis
an bhfómhar. Ach fad is a bhíonn na daoine bochta ag glacadh leis na sean-uirlisí
mar a bhíodh --- \VUgras{speala, corráin} agus \VUgras{súistí}
\textsuperscript{(40)} --- baineann na húinéirí saibhre a gcruithneacht nó a
\VUgras{seagal} \textsuperscript{(43)} leis an \VUgras{inneall bainte}
\textsuperscript{(42)}. Ansin, gan na punanna a thabhairt chuig an scioból,
déanann siad \VUgras{cruacha} móra \textsuperscript{(44)} ar an láthair féin.

%%K t08-c1-11-1 p
11. --- Ansin tugtar an \VUgras{inneall buailte} \textsuperscript{(47)}, a
chasann an \VUgras{lócamóibíl} \textsuperscript{(45)} tríd an \VUgras{gcrios
tarchurtha} \textsuperscript{(46)}, agus mar sin scartar an gráinne ón
\VUgras{tuí} gan an ghort ina bhfuil sé curtha a fhágáil.

%%K t08-tit-4 sub
\VUsaut{5.0mm}
\VUpk{10.2pt}{40}{An Stoirm.}
\VUsaut{3.4mm}

%%K t08-c1-12-1 p
12. --- Le linn an fhómhair briseann \VUgras{stoirmeacha} láidre amach go minic
\textsuperscript{(53)}. Ar dtús cluintear i bhfad uainn tormán na
\VUgras{toirní}; éiríonn \VUgras{gaoth aniar} láidir \textsuperscript{(54)} agus
lúbann sí go dtí an cnead na crainn is mó sa \VUgras{choill}
\textsuperscript{(49)}. Gabhann \VUgras{scamaill} dhubha \textsuperscript{(56)} an
spéir; trasnaíonn zigeanna dallta na \VUgras{tintrí} \textsuperscript{(55)} iad.
Tosaíonn an \VUgras{bháisteach} ag titim, agus uaireanta an \VUgras{chloch
shneachta} freisin, ag leagan an arbhair atá réidh don fhómhar.

%%K t08-c1-13-1 p
13. --- Is minic a lasann an tintreach foirgneamh éigin. Mar sin, anuraidh dódh
díon an \VUgras{mhuilinn ghaoithe} \textsuperscht{(50)} go talamh, agus briseadh a
dhá \VUgras{sciathán} \textsuperscript{(51)}.

%%K t08-c1-14-1 p
14. --- Tar éis cúpla uair an chloig filleann an suaimhneas. Ag bun na spéire
nochtann \VUgras{tuar ceatha} lonrach \textsuperscript{(52)}, agus tosaíonn an
nádúr arís ar an mbeatha a bhí briste ar feadh tamaill. Filleann muintir an
tsráidbhaile, a bhí ar foscadh sna \VUgras{botháin} \textsuperscript{(48)}, ar an
obair agus tosaíonn siad go cróga ar an damáiste a rinneadh a dheisiú.

%%K t08-tit-5 sub
\VUsaut{6.0mm}
\VUcentre{13.2pt}{90}{\textit{An Dara Radharc.}}
\VUsaut{3.6mm}

%%K t08-tit-6 sub
\VUpk{10.2pt}{40}{An tSealgaireacht. --- Baint na bhFíonchaor.}
\VUsaut{1.4mm}
\VUpk{10.2pt}{40}{An Iascaireacht.}
\VUsaut{4.0mm}

%%K t08-c2-01-1 p
1. --- Faoi dheireadh mhí \VUgras{Lúnasa} nó faoi lár mhí \VUgras{Mheán Fómhair},
de réir an cheantair, osclaítear an tsealgaireacht. Chomh luath agus a mheallann
céad ghathanna na gréine na \VUgras{laghairteanna} \textsuperscript{(2)} as a
bpoill, bíonn an \VUgras{sealgaire} \textsuperscript{(5)} ag dul amach lena
\VUgras{mhadraí} \textsuperscript{(4)}.

%%K t08-c2-02-1 p
2. --- Chuir sé air a \VUgras{chulaith sealgaireachta} láidir \textsuperscript{(9)}
d'éadach tiubh agus na \VUgras{loirgníní} leathair, a ligfidh dó siúl tríd an
bhféar fliuch, mar a bhfuil na \VUgras{buafanna} \textsuperscript{(1)} i bhfolach;
thóg sé an \VUgras{gunna} \textsuperscript{(6)} agus an \VUgras{mála sealgaireachta}
\textsuperscript{(10)} --- agus chun bóthair.

%%K t08-c2-03-1 p
3. --- Tugann díograis na sealgaireachta é go lár an fhíonghoirt féin, mar a bhfuil
an bhaint faoi lánseol. Ligeann páistí an fhíonadóra, a fhaireann na gabhair in
aice leis an mbóthar de ghnáth, dóibh inniu innilt mar is mian leo agus, tar éis
dóibh an sean-\VUgras{phocán} \textsuperscript{(3)} a cheangal de chuaille --- pocán
a bhainfeadh leas as an tsaoirse cinnte chun teitheadh --- glacann siadsan a gcuid
den ghreann freisin. Tógann duine acu triopall as an \VUgras{gciseán fíonchaor}
\textsuperscript{(12)} agus bíonn sé ag spraoi ag fáisceadh an \VUgras{tsú}
\textsuperscript{(14)} isteach i \VUgras{gcupán} \textsuperscript{(13)}, chun go
ndéanfar múst (fíon gan choipeadh) de. Corónaíonn duine eile é féin le
\VUgras{fleasc duilleog} \textsuperscript{(15)}, atá díreach fite aige. Lena
dtaobh tagann \VUgras{gealbhain} dhána \textsuperscript{(16)} go dtí an fáisceán
féin, chun fíonchaora a ghoid.

%%K t08-c2-04-1 p
4. --- Fad is atá na páistí ag spraoi, bíonn na fir ag obair. Iompraíonn na
\VUgras{buanaithe} \textsuperscript{(18)} na fíonchaora go dtí an siléar i
\VUgras{gciseáin dhroma} \textsuperscript{(17)}; deisíonn an \VUgras{cúipéir}
\textsuperscript{(19)} na \VUgras{bairillí} \textsuperscript{(20)}, féachann sé an
bhfuil na \VUgras{cláracha} \textsuperscript{(22)} agus na \VUgras{tóin}
\textsuperscript{(21)} slán, agus athraíonn sé na \VUgras{fáinní} briste
\textsuperscript{(23)}. Is é freisin a rinne na \VUgras{dabhcha} móra
\textsuperscript{(25)}, ina ndoirtear na \VUgras{fíonchaora}
\textsuperscript{(26)}, chun go gcoipfidís.

%%K t08-c2-05-1 p
5. --- Nuair a bheidh an choipeadh thart, doirtear an fíon, agus cuirtear an
fuílleach faoin \VUgras{bhfáisceán} \textsuperscript{(28)}; ag casadh an
\VUgras{scriú} mhóir go mall \textsuperscript{(29)}, fáisctear an sú, a ritheann
isteach sa \VUgras{dabhach} \textsuperscript{(32)} agus a thugann \VUgras{fíon}
\textsuperscript{(31)} fós.

%%K t08-tit-8 sub
\VUsaut{6.0mm}
\VUpk{10.2pt}{40}{An Beoir agus an Ceirtlis.}
\VUsaut{3.4mm}

%%K t08-c2-06-1 p
6. --- Ar bhalla an \VUgras{tsiléir} \textsuperscript{(27)} casann craobh
\VUgras{leannlusa} \textsuperscript{(30)}. Anseo níl ann ach maisiú, ach i dtíortha
eile, mar a bhfuil an fhíniúin an-ghann, saothraítear an leannlus chun
\VUgras{beoir} a dhéanamh.

%%K t08-c2-07-1 p
7. --- Cromann an \VUgras{crann úll} beag \textsuperscript{(33)} faoi na húlla, a
thabharfaidh \VUgras{ceirtlis} den scoth nuair a aibeoidh siad. Ina
\VUgras{gcliabhán} \textsuperscript{(34)} féachann an \VUgras{canáraí}
\textsuperscript{(35)} agus an \VUgras{lasair choille} \textsuperscript{(36)} le
héad ar na gealbhain atá ag spraoi faoi shaoirse.

%%K t08-c2-08-1 p
8. --- Chomh fada agus a shroicheann an tsúil, ar fhána na gcnoc tá sraitheanna
\VUgras{cuaillí} \textsuperscript{(40)}, a thacaíonn leis na \VUgras{fíniúna}
iontacha \textsuperscript{(39)}, trom le \VUgras{triopaill}.

%%K t08-c2-08-2 p
Ar feadh na bliana ar fad shaothraigh an \VUgras{fíonadóir}
\textsuperscript{(42)} chun aire a thabhairt dá \VUgras{fhíonghort}
\textsuperscript{(38)}, agus seo é luaite as a chuid oibre le fómhar maith.
Líonann na \VUgras{bainteoirí} \textsuperscript{(41)} na dabhcha atá curtha ar
chairt, atá á tarraingt ag \VUgras{miúileanna} \textsuperscript{(43)}, agus tá gach
duine meidhreach.

%%K t08-tit-9 sub
\VUsaut{5.0mm}
\VUpk{10.2pt}{40}{An Iascaireacht.}
\VUsaut{3.4mm}

%%K t08-c2-09-1 p
9. --- Mar an gcéanna do na \VUgras{hiascairí slaite} \textsuperscript{(44)}, a
tháinig ag breacadh an lae chun suí cois uisce lena \VUgras{slata}
\textsuperscript{(45)}.

%%K t08-c2-09-2 p
Beireann an \VUgras{t-iasc} \textsuperscript{(47)} ar an duán chomh luath agus a
ligeann siad a \VUgras{ndorú} \textsuperscript{(46)} san uisce, agus is ar éigean a
bhíonn am acu an \VUgras{bhaoite} a athrú sula mbíonn íospartach nua ag preabadh ag
ceann an dorú.

%%K t08-c2-09-3 p
Mar sin féin beireann an \VUgras{t-iascaire eangaí} \textsuperscript{(48)}, a
chaitheann a líon ón \VUgras{mbád} \textsuperscript{(49)} go lár na \VUgras{habhann}
\textsuperscript{(50)}, i bhfad níos mó.

%%K t08-c2-10-1 p
10. --- Is é an fómhar am na siúlóidí breátha: de shiúl na gcos, ar
\VUgras{muin capaill} \textsuperscript{(52)}, ar \VUgras{rothar}
\textsuperscript{(53)}, i \VUgras{ngluaisteán} \textsuperscript{(54)}. Tá na
\VUgras{bóithre} \textsuperscript{(51)} glan, agus baineann na daoine leas as na
laethanta breátha deireanacha, mar tá na \VUgras{fáinleoga}
\textsuperscript{(56)} ag cruinniú ar na díonta cheana chun eitilt chuig na tíortha
teo. Tá duilleoga na \VUgras{gcoll} sean \textsuperscript{(57)} ag buíochan, tá na
\VUgras{cnónna} ag titim, agus na \VUgras{crainn} arda, atá ina
\VUgras{ndoire} \textsuperscript{(58)} ar an \VUgras{ard}
\textsuperscript{(59)}, beidh siad lom gan mhoill.

%%K t08-c2-11-1 p
11. --- Éiríonn an \VUgras{ghrian} \textsuperscript{(60)} níos déanaí, éiríonn na
laethanta níos giorra, tosaíonn fuacht sách géar le mothú ar \VUgras{maidin}, agus
tá scataí \VUgras{creabhar} \textsuperscript{(61)} ag fágáil ár dtíre
neamhfháiltiúla cheana.
\VUsaut{6.0mm}
\VUfilet{22mm}
